\documentclass[11pt]{article}
% ======================================================================
%  examstyle.tex — EPFL-style exam layout for the weekly Time Series exams
%  Shared by every Exam_Week_NN(.tex / _Solutions.tex). Compiles with pdflatex.
% ======================================================================
\usepackage[utf8]{inputenc}
\usepackage[T1]{fontenc}
\usepackage[english]{babel}
\usepackage[a4paper,margin=2.4cm]{geometry}
\usepackage{amsmath,amssymb,amsthm,mathtools}
\usepackage{bm}
\usepackage{enumitem}
\usepackage{array}

\setlength{\parindent}{0pt}
\setlength{\parskip}{0.5em}
\emergencystretch=3em
\allowdisplaybreaks[1]

% ---------- math macros (same names as the rest of the project) ----------
\newcommand{\E}{\mathbb{E}}
\DeclareMathOperator{\Var}{Var}
\DeclareMathOperator{\Cov}{Cov}
\DeclareMathOperator{\Corr}{Corr}
\DeclareMathOperator*{\argmin}{arg\,min}
\DeclareMathOperator{\MSE}{MSE}
\DeclareMathOperator{\trace}{tr}
\newcommand{\Reals}{\mathbb{R}}
\newcommand{\Ints}{\mathbb{Z}}
\newcommand{\Comp}{\mathbb{C}}
\newcommand{\Nat}{\mathbb{N}}
\newcommand{\Prob}{\mathbb{P}}
\newcommand{\Tr}{^{\mathsf{T}}}
\newcommand{\conj}[1]{\overline{#1}}
\newcommand{\abs}[1]{\left\lvert #1 \right\rvert}
\newcommand{\norm}[1]{\left\lVert #1 \right\rVert}
\newcommand{\ind}{\mathbf{1}}
\newcommand{\eps}{\varepsilon}
\newcommand{\diff}{\nabla}
\newcommand{\acvs}{\gamma}
\newcommand{\acf}{\rho}
\newcommand{\pacf}{\alpha}
\newcommand{\sdf}{\mathcal{S}}
\newcommand{\Bsh}{B}
\newcommand{\iid}{\overset{\text{iid}}{\sim}}

\setlist[enumerate]{leftmargin=2em,topsep=2pt,itemsep=4pt}
\setlist[itemize]{leftmargin=1.6em,topsep=2pt,itemsep=2pt}

\newcounter{exo}

% ---------- exam cover header :  \examheader{exam title}{duration} ----------
\newcommand{\examheader}[2]{%
  \thispagestyle{empty}
  \begin{center}
    {\Large\bfseries Time Series}\\[3pt]
    Spring Semester 2026, EPFL\\
    \'Ecole Polytechnique F\'ed\'erale de Lausanne\\
    Prof.\ Dr.\ Sofia C.\ Olhede\\[7pt]
    \hrule height 0.8pt
    \vspace{9pt}
    {\large\bfseries Practice Exam --- #1}\\[4pt]
    {\normalsize Duration: #2 \quad$\cdot$\quad No calculator \quad$\cdot$\quad Answer on paper}\\
    \vspace{8pt}
    \hrule height 0.8pt
  \end{center}
  \vspace{14pt}
  \begin{tabular}{@{}p{8.2cm}@{}p{8cm}@{}}
    Name:\enspace\hrulefill & Surname:\enspace\hrulefill
  \end{tabular}\\[14pt]
  SCIPER (Student Card Number):\enspace\hrulefill
  \vspace{16pt}
}

% ---------- points table (fixed: 4 problems) ----------
\newcommand{\pointstable}{%
  \begin{center}
  \renewcommand{\arraystretch}{1.5}
  \begin{tabular}{|p{5cm}|p{4cm}|}
    \hline
    \textbf{Exercise} & \textbf{Points}\\\hline
    1 & \\\hline
    2 & \\\hline
    3 & \\\hline
    4 & \\\hline
    \textbf{Total} & \\\hline
  \end{tabular}
  \end{center}
}

% ---------- problem heading (exam: one problem per page) ----------
\newcommand{\problem}[1]{%
  \clearpage
  \stepcounter{exo}%
  \begin{center}\textbf{\large Problem \theexo\ (#1 points)}\end{center}
  \vspace{4pt}
}
% ---------- answer space: fill the statement page + one blank answer page ----------
% Put \answerpage at the END of each problem so every exercise gets its own page(s).
\newcommand{\answerpage}{\par\vfill\newpage\null\newpage}

% ---------- solutions document ----------
\newcommand{\solheader}[1]{%
  \begin{center}
    {\Large\bfseries Time Series --- Solutions}\\[3pt]
    {\large Practice Exam --- #1}\\[2pt]
    EPFL, MATH-342 \quad$\cdot$\quad Prof.\ S.\ C.\ Olhede\\[5pt]
    \hrule height 0.8pt
  \end{center}
  \vspace{10pt}
}
\newcommand{\solproblem}[1]{%
  \bigskip
  \stepcounter{exo}%
  {\large\bfseries Problem \theexo\ (#1 points)}\par\nobreak\vspace{2pt}
}
% Rappel de l'énoncé avant la solution (dans les corrigés)
\newcommand{\statementstart}{\par\vspace{1pt}{\bfseries\itshape Statement.}\par\nobreak\begingroup\small\itshape}
\newcommand{\statementend}{\par\endgroup\vspace{5pt}\hrule\vspace{6pt}{\bfseries Solution.}\par\nobreak\vspace{2pt}}

\begin{document}

\examheader{Week 2: ARMA Models and the PACF}{60 minutes}
\pointstable
\begin{center}\itshape Justify every answer. Throughout, $\eps_t$ is a mean-zero white noise of variance $\sigma^2$.\end{center}

% ---------------------------------------------------------------
\problem{5}
\textbf{(a)} State the definition of white noise and write down its ACVS $\acvs_\tau$ and ACF $\acf_\tau$. Explain in one sentence why ``uncorrelated'' is weaker than ``independent''. \hfill(1.5 pts)

\textbf{(b)} Consider the MA(2) process
\[
X_t=\eps_t-\theta_1\eps_{t-1}-\theta_2\eps_{t-2}.
\]
Compute the mean and the autocovariance $\acvs_\tau$ for \emph{all} $\tau\in\Ints$, and write down $\acf_1$ and $\acf_2$. \hfill(2.5 pts)

\textbf{(c)} For which lags $\tau$ does the ACVS of an MA($q$) vanish? State the general rule and explain how it is used to read off the order $q$ from an estimated ACF. \hfill(1 pt)
\answerpage

% ---------------------------------------------------------------
\problem{5}
Consider the AR(1) process $X_t=\phi X_{t-1}+\eps_t$.

\textbf{(a)} State the condition on $\phi$ for the process to be causal, and explain it in terms of the root of the AR polynomial $\Phi(z)=1-\phi z$. \hfill(1 pt)

\textbf{(b)} Assuming $\abs\phi<1$, derive the causal MA($\infty$) representation $X_t=\sum_{k\ge0}\psi_k\eps_{t-k}$ by iterating the recursion. Give the $\psi$-weights and verify $\sum_k\abs{\psi_k}<\infty$. \hfill(2 pts)

\textbf{(c)} Using the representation in (b), show that
\[
\acvs_\tau=\frac{\sigma^2}{1-\phi^2}\,\phi^{\abs\tau},
\qquad
\acf_\tau=\phi^{\abs\tau}.
\]
Why does this ACF \emph{tail off} rather than cut off, and what does that say about identifying $p$ from the ACF? \hfill(2 pts)
\answerpage

% ---------------------------------------------------------------
\problem{5}
\textbf{(a)} Consider the AR(2) process $X_t=\phi_1 X_{t-1}+\phi_2 X_{t-2}+\eps_t$ with $\phi_1=\tfrac12,\ \phi_2=-\tfrac14$. Using the causality conditions $\phi_1+\phi_2<1,\ \phi_2-\phi_1<1,\ \abs{\phi_2}<1$, verify that the process is causal. \hfill(1 pt)

\textbf{(b)} From the Yule--Walker equations $\acf_\tau=\phi_1\acf_{\tau-1}+\phi_2\acf_{\tau-2}$ (with $\acf_0=1,\ \acf_{-1}=\acf_1$), compute $\acf_1,\acf_2,\acf_3$. \hfill(2 pts)

\textbf{(c)} Using the Durbin--Levinson recursion
\[
\pacf_{1,1}=\acf_1,\qquad
\pacf_{\tau,\tau}=\frac{\acf_\tau-\sum_{j=1}^{\tau-1}\pacf_{\tau-1,j}\,\acf_{\tau-j}}{1-\sum_{j=1}^{\tau-1}\pacf_{\tau-1,j}\,\acf_{j}},
\qquad
\pacf_{\tau,j}=\pacf_{\tau-1,j}-\pacf_{\tau,\tau}\,\pacf_{\tau-1,\tau-j},
\]
compute $\pacf_1,\pacf_2,\pacf_3$ from the values in (b). Read off the implied model order and state the ACF/PACF identification rule it illustrates. \hfill(2 pts)
\answerpage

% ---------------------------------------------------------------
\problem{5}
Consider the ARMA(1,1) process $X_t=\phi X_{t-1}+\eps_t+\vartheta\eps_{t-1}$ with $\phi=\tfrac12,\ \vartheta=\tfrac14$.

\textbf{(a)} Write the model as $\Phi(\Bsh)X_t=\Theta(\Bsh)\eps_t$ and check causality and invertibility from the roots of $\Phi(z)$ and $\Theta(z)$. \hfill(1.5 pts)

\textbf{(b)} Derive the $\psi$-weights of the causal MA($\infty$) representation $X_t=\sum_{j\ge0}\psi_j\eps_{t-j}$ by matching coefficients in $\Phi(\Bsh)\psi(\Bsh)=\Theta(\Bsh)$. Give $\psi_0,\psi_1,\psi_2$ and the general formula. \hfill(2 pts)

\textbf{(c)} \emph{(Week 1 revision.)} Recall that a candidate sequence $\{\acvs_\tau\}$ is a valid ACVS only if it is positive semi-definite. Decide whether
\[
\acvs_0=1,\qquad \acvs_{\pm1}=0.7,\qquad \acvs_\tau=0\ (\abs\tau\ge2)
\]
is a valid ACVS of some stationary process, by computing the spectral density $\sdf(f)=\sum_\tau\acvs_\tau e^{-2\pi i f\tau}$ and checking its sign. \hfill(1.5 pts)
\answerpage

\end{document}
